%----------------------------------------------------------------------------------------
%    PACKAGES AND THEMES
%----------------------------------------------------------------------------------------

\documentclass[aspectratio=169,xcolor=dvipsnames]{beamer}
\usetheme{SimplePlus}
\setbeamercolor{normal text}{fg=black}
\setbeamercolor{frametitle}{fg=black}

\usepackage{hyperref}
\usepackage{graphicx}
\usepackage{booktabs}
\usepackage[utf8]{inputenc}
\usepackage[spanish, provide=*]{babel}
\usepackage{pgfplots}
\pgfplotsset{compat=1.17}

\useinnertheme{circles}

%----------------------------------------------------------------------------------------
%    TITLE PAGE
%----------------------------------------------------------------------------------------

\title[]{La Función Drop-Wave}
\author[Enrique Flores]{Enrique Lisandro Flores Brito}

\institute
{Universidad Autónoma del Estado de Morelos\\
    Instituto de Investigación en Ciencias Básicas y Aplicadas\\
    Centro de Investigación en Ciencias\\
    Lic. en Inteligencia Artificial\\[0.5cm] % Espacio para separar
    Búsqueda de Soluciones e Inferencia Bayesiana\\
    Profesor: Dr. Jorge Hermosillo Valadez}
\date{ }

\usepackage{tikz}
\usetikzlibrary{shapes.arrows}
\usepackage{tcolorbox}
%\usepackage{palatino}
\setbeamercolor{normal text}{fg=black}
\usefonttheme{default} % Typeset using the default sans serif font
\setbeamercolor{frametitle}{fg=black}
\setbeamercolor{section in toc}{fg=black}

%	DEFINIR ENUNCIADOS IMPORTANTES
%\DeclarePairedDelimiterX\set[1]\lbrace\rbrace{\def\given{\;\delimsize\vert\;}#1}
\newtheorem{teo}{Teorema}[section] %	Declarar Teoremas con la numeración de la sección
\newtheorem{coro}{Corolario}[teo]  %	Declarar Corolarios con la numeración de los Teoremas
\newtheorem*{coro*}{Corolario}
\newtheorem{lema}{Lema}[section]   %	Seguir la misma numeración de los Teoremas
\newtheorem*{prop*}{Proposición}
\newtheorem{prop}{Proposición}

\theoremstyle{definition}
\newtheorem{defi}{Definición}[section]

\theoremstyle{remark}
\newtheorem{obs}{Observación}

%	Comandos personalizados

\newcommand{\ZZ}{\mathbb{Z}}	% Así nos ahorramos escribir todo el comando para poner el símbolo de los enteros, reales, etc...
\newcommand{\NN}{\mathbb{N}}
\newcommand{\RR}{\mathbb{R}}

%----------------------------------------------------------------------------------------
%    PRESENTATION SLIDES
%---------------------------------------------------------------------

\begin{document}
    \frame{\titlepage}

    \begin{frame}{Introducción: Función Drop-Wave}
    \begin{columns}[c] % Columnas centradas verticalmente

        % Columna Izquierda: Información esencial
        \column{0.45\textwidth}
        \begin{block}{Características Clave}
            \begin{itemize}
                \item \textbf{Tipo:} Función de prueba (benchmark) multimodal.
                \item \textbf{Desafío:} Múltiples mínimos locales en forma de ondas concéntricas.
                \item \textbf{Mínimo Global:}
                \[ f(0,0) = -1 \]
            \end{itemize}
        \end{block}

        \begin{tcolorbox}[colback=blue!5,colframe=blue!75!black,title=Ecuación]
            \small
            \[ f(\mathbf{x}) = - \frac{1 + \cos(12 \sqrt{x_1^2 + x_2^2})}{0.5(x_1^2 + x_2^2) + 2} \]
        \end{tcolorbox}

        % Columna Derecha: Gráfico en LaTeX
        \column{0.55\textwidth}
        \begin{tikzpicture}
            \begin{axis}[
                width=\textwidth,
                view={-35}{25},
                enlargelimits=false,
                grid=major,
                domain=-5.12:5.12,
                y domain=-5.12:5.12,
                samples=40, % Resolución balanceada para velocidad de compilación
                colormap/jet,
                colorbar,
                colorbar style={width=0.2cm},
                xlabel=\tiny $x_1$,
                ylabel=\tiny $x_2$,
                zlabel=\tiny $f$,
                zmin=-1.1, zmax=0
            ]
                \addplot3[
                    surf,
                    shader=interp,
                ]
                {-(1 + cos(deg(12 * sqrt(x^2 + y^2)))) / (0.5 * (x^2 + y^2) + 2)};
            \end{axis}
        \end{tikzpicture}
    \end{columns}
    \end{frame}

    \begin{frame}{Optimización con Algoritmos Genéticos}
        \begin{block}{Representación del Problema}
            Para la función Drop-Wave, el enfoque de \textbf{Codificación Real} es el más eficiente:
            \begin{itemize}
                \item \textbf{Individuo (Cromosoma):} Un vector $\vec{x} = (x_1, x_2)$ donde cada gen es un número flotante.
                \item \textbf{Espacio de Búsqueda:} Definido estrictamente en el intervalo $x_i \in [-5.12, 5.12]$.
                \item \textbf{Población Inicial:} Generada aleatoriamente mediante una distribución uniforme para asegurar cobertura total del dominio.
            \end{itemize}
        \end{block}

        \begin{block}{Función de Aptitud (Fitness)}
            Dado que buscamos el mínimo global, la función de aptitud debe ser la misma función pero negativa: 
            \begin{itemize}
                \item \textbf{Criterio:} Minimizar $f(\vec{x})$ es equivalente a maximizar $\Phi(\vec{x}) = -f(\vec{x})$.
                \item \textbf{Evaluación:} Se calcula la aptitud de cada candidato mediante la función Drop-Wave para cuantificar su ventaja adaptativa.
            \end{itemize}
        \end{block}
    \end{frame}

    \begin{frame}[t]{Optimización con Algoritmos Genéticos}
    \vspace{0.2cm}

    \begin{columns}[t]
        \column{0.5\textwidth}
        \begin{block}{Operadores Evolutivos}
            \begin{itemize}
                \item \textbf{Selección:} Por torneo.
                \item \textbf{Cruza:} Recombinación aritmética.
                \item \textbf{Mutación:} Ruido Gaussiano para saltar crestas.
            \end{itemize}
        \end{block}

        \column{0.5\textwidth}
        \begin{tcolorbox}[colback=red!5, colframe=red!75!black, title=Desafío, top=1mm, bottom=1mm]
            \small
            La \textbf{multimodalidad} castiga la falta de diversidad. Sin mutación suficiente, la población converge en anillos exteriores, perdiendo el pozo central en $(0,0)$.
        \end{tcolorbox}
    \end{columns}

    \vspace{0.3cm}

    \begin{block}{Condición de Parada}
        Convergencia en región mínima o alcance del límite de generaciones, buscando el valor teórico de $-1$.
    \end{block}
\end{frame}

\begin{frame}
    \frametitle{Referencias}

    \begin{thebibliography}{9}
        \bibitem[Molga and Smutnicki, 2005]{molga2005}
        M. Molga and C. Smutnicki (2005)
        \newblock \emph{Test functions for optimization needs}
        \newblock \emph{Test functions for optimization needs}, 101(48), 32.

        \bibitem[Surjanovic and Bingham, 2013]{surjanovic2013}
        S. Surjanovic and D. Bingham (2013)
        \newblock \emph{Virtual Library of Simulation Experiments: Drop-Wave Function}
        \newblock Simon Fraser University. Retrieved from \url{https://www.sfu.ca/~ssurjano/drop.html}
    \end{thebibliography}
\end{frame}
\end{document}
